\begin{abstract}
EIP-7702 lets an externally owned account delegate execution to contract code, but authorizing an untrusted delegate can expose the account before any harmful transaction history exists. We study whether a defender can screen the delegate at this pre-authorization point using runtime bytecode alone. AuthGuard-7702 is an evaluation-grade research prototype that deterministically extracts bytecode features and returns a risk score without decompilation. We construct a task-aligned dataset by applying an outcome-blind policy to delegation designators and conflicting exact-bytecode labels, then use family-grouped evaluation to hold related bytecode together. Against USENIX-artifact positives and rule-silent weak negatives, AuthGuard attains family-grouped and seeded-random AUPRC of $0.881\pm0.028$ and $0.975\pm0.012$, respectively. Under the stricter augmentation protocol, held-out pure-M0 F200 fold-mean recall/AUPRC/FPR changes from $0.484/0.561/0.217$ to $0.727/0.758/0.174$; the family-clustered pooled recall difference is $+0.253$ (95\% CI $[0.144,0.379]$). These labels do not constitute perfect malicious/benign ground truth, and the prototype excludes wallet, network, authorization-parsing, and warning-interface integration. The results support bytecode scoring as a complementary warning signal while leaving broader obfuscation and compound transformations unresolved.
\end{abstract}

\begin{IEEEkeywords}
EIP-7702, smart-contract security, bytecode classification, family-grouped evaluation, adversarial robustness, risk scoring
\end{IEEEkeywords}
